\section{Metric Spaces \& General Topology}

\subsection{Metric Spaces, Topological Properties, and Completeness}

We now abstract the topological concepts developed for $\mathbb{R}$ in Chapter 1 to arbitrary sets equipped with a notion of distance. This abstraction provides the theoretical foundation for functional analysis, where the points in the space are often themselves functions.

\begin{definition}
A \textbf{metric space} is a pair $(X, d)$, where $X$ is a set and $d: X \times X \to \mathbb{R}$ is a function (the metric) satisfying the following properties for all $x, y, z \in X$:
\begin{enumerate}
    \item \textbf{Positivity:} $d(x, y) \ge 0$, and $d(x, y) = 0$ if and only if $x = y$.
    \item \textbf{Symmetry:} $d(x, y) = d(y, x)$.
    \item \textbf{Triangle Inequality:} $d(x, z) \le d(x, y) + d(y, z)$.
\end{enumerate}
\end{definition}

The topology of a metric space is generated by open balls. For $x \in X$ and $r > 0$, the open ball of radius $r$ centered at $x$ is defined as $B_r(x) = \{ y \in X \mid d(x, y) < r \}$. A subset $U \subset X$ is \textbf{open} if for every $x \in U$, there exists an $r > 0$ such that $B_r(x) \subset U$. A subset $F \subset X$ is \textbf{closed} if its complement $X \setminus F$ is open.

\begin{definition}
A sequence $(x_n)_{n=1}^\infty$ in a metric space $(X, d)$ is a \textbf{Cauchy sequence} if for every $\epsilon > 0$, there exists an $N \in \mathbb{N}$ such that for all $m, n \ge N$, $d(x_m, x_n) < \epsilon$. The metric space $(X, d)$ is \textbf{complete} if every Cauchy sequence in $X$ converges to a limit in $X$. 
\end{definition}

A particularly important class of complete metric spaces arises from vector spaces equipped with a norm.

\begin{definition}
A \textbf{Banach space} is a vector space $V$ over $\mathbb{R}$ or $\mathbb{C}$ equipped with a norm $\|\cdot\|$ such that the metric space induced by the metric $d(u, v) = \|u - v\|$ is complete.
\end{definition}

Completeness is a powerful property because it guarantees the existence of a limit based solely on the internal behavior of the sequence, without prior knowledge of the limit point. This principle is exploited in one of the most fundamental results in non-linear analysis: the Banach Fixed-Point Theorem.

\begin{definition}
Let $(X, d)$ be a metric space. A map $T: X \to X$ is a \textbf{contraction mapping} if there exists a constant $c \in [0, 1)$ such that for all $x, y \in X$:
$$ d(T(x), T(y)) \le c \, d(x, y) $$
\end{definition}

\begin{theorem}[Banach Fixed-Point Theorem]
Let $(X, d)$ be a non-empty complete metric space, and let $T: X \to X$ be a contraction mapping with contraction constant $c \in [0, 1)$. Then $T$ admits a unique fixed point $x^* \in X$ (i.e., $T(x^*) = x^*$).
\end{theorem}

\begin{proof}
We first prove uniqueness. Suppose $x, y \in X$ are both fixed points of $T$, such that $T(x) = x$ and $T(y) = y$. Applying the contraction property yields:
$$ d(x, y) = d(T(x), T(y)) \le c \, d(x, y) $$
This implies $(1 - c) d(x, y) \le 0$. Since $c < 1$, we have $1 - c > 0$. Because distance is non-negative ($d(x, y) \ge 0$), the only way this inequality can hold is if $d(x, y) = 0$. By the positivity axiom of a metric space, $x = y$. Thus, if a fixed point exists, it is strictly unique.

To prove existence, we construct a sequence via Picard iteration. Choose an arbitrary initial point $x_0 \in X$. Define the sequence $(x_n)_{n=0}^\infty$ recursively by $x_{n} = T(x_{n-1})$ for all $n \ge 1$. We must show that $(x_n)$ is a Cauchy sequence.

Consider the distance between consecutive terms. By repeated application of the contraction property:
$$ d(x_{n+1}, x_n) = d(T(x_n), T(x_{n-1})) \le c \, d(x_n, x_{n-1}) \le \dots \le c^n d(x_1, x_0) $$
Let $m > n \ge 0$. We apply the triangle inequality repeatedly to bound the distance between $x_m$ and $x_n$:
$$ d(x_m, x_n) \le d(x_m, x_{m-1}) + d(x_{m-1}, x_{m-2}) + \dots + d(x_{n+1}, x_n) $$
Substituting the bound for consecutive terms, we obtain:
$$ d(x_m, x_n) \le (c^{m-1} + c^{m-2} + \dots + c^n) d(x_1, x_0) = c^n \sum_{k=0}^{m-n-1} c^k \, d(x_1, x_0) $$
Because $c \in [0, 1)$, the infinite geometric series $\sum_{k=0}^\infty c^k$ converges to $\frac{1}{1-c}$. Therefore, we can strictly bound the finite sum:
$$ d(x_m, x_n) \le c^n \frac{1}{1 - c} d(x_1, x_0) $$
Let $\epsilon > 0$. Since $c < 1$, $\lim_{n \to \infty} c^n = 0$. Thus, there exists an $N \in \mathbb{N}$ such that for all $n \ge N$:
$$ c^n < \frac{\epsilon (1 - c)}{d(x_1, x_0)} $$
(If $d(x_1, x_0) = 0$, the sequence is constant and trivially Cauchy). For all $m > n \ge N$, this forces $d(x_m, x_n) < \epsilon$. Thus, $(x_n)$ is a Cauchy sequence.

Because $X$ is a complete metric space, the Cauchy sequence $(x_n)$ converges to a limit $x^* \in X$. We now verify that $x^*$ is a fixed point. A contraction mapping is globally uniformly continuous (letting $\delta = \epsilon$ works for any $c < 1$). Because $T$ is continuous, the limit commutes with the function application:
$$ T(x^*) = T\left(\lim_{n \to \infty} x_n\right) = \lim_{n \to \infty} T(x_n) = \lim_{n \to \infty} x_{n+1} = x^* $$
Hence, $x^*$ is a fixed point of $T$, completing the proof.
\end{proof}
